\bibitem{Kontsevich}
M. Kontsevich, \emph{Deformation quantization of Poisson manifolds},
Lett. Math. Phys. \textbf{66} (2003), 157--216.
\url{https://arxiv.org/abs/q-alg/9709040}.
Section~4.6.2のHKR比較を参照する。

\bibitem{CT}
A. C\u{a}ld\u{a}raru and J. Tu,
\emph{Curved $A_\infty$ algebras and Landau--Ginzburg models},
New York J. Math. \textbf{19} (2013), 305--342.
Section~2, Remark~2.8, Theorem~2.12.
\url{https://nyjm.albany.edu/j/2013/19-16v.pdf}.

\bibitem{PMayanskiy}
E. Mayanskiy, \emph{Poisson cohomology of two Fano threefolds},
arXiv:1303.5916, Theorem~1およびSection~2の係数計算。
\url{https://arxiv.org/abs/1303.5916}.
同Theorem~1はLoray--Pereira--Touzetの既知結果として述べられている。

\bibitem{JLLZ}
A. Jacovskis, X. Lin, Z. Liu and S. Zhang,
\emph{Infinitesimal categorical Torelli theorems for Fano threefolds},
J. Pure Appl. Algebra \textbf{227} (2023), 107418.
\url{https://doi.org/10.1016/j.jpaa.2023.107418}.
著者版（2022年9月25日）のTheorem~3.10, Lemma~3.13およびTheorem~4.4(2)
（印刷pp.~9--11, 15）を参照する。
\url{https://ajacovskis.github.io/hochschild_only_2.pdf}.

\bibitem{AMChen}
X.-W. Chen, \emph{Module factorizations}, arXiv:2402.11613v3 (2025).
Theorem~2.10, Corollary~2.11, Theorem~7.2。
\url{https://arxiv.org/html/2402.11613v3}.
余核の三角同値を使用し、enhancementは本文でexact圏とdg商により構成する。

\bibitem{AZ}
M. Artin and J. J. Zhang, \emph{Noncommutative projective schemes},
Adv. Math. \textbf{109} (1994), 228--287.
\url{https://doi.org/10.1006/aima.1994.1087}.

\bibitem{Orlov}
D. Orlov, \emph{Derived categories of coherent sheaves and triangulated categories of singularities},
arXiv:math/0503632v3 (2009), Theorem~2.5。
\url{https://arxiv.org/abs/math/0503632}.

\bibitem{BS}
J. Burke and G. Stevenson, \emph{The derived category of a graded Gorenstein ring},
arXiv:1507.00830v1 (2015), Theorem~6.4およびその証明。
\url{https://arxiv.org/abs/1507.00830}.
付録では剰余体上のwindowとVerdier商の向きの同定に用いる。

\bibitem{PFCRVdB}
D. Calaque, C. A. Rossi and M. Van den Bergh,
\emph{C\u{a}ld\u{a}raru's conjecture and Tsygan's formality},
Ann. of Math. (2) \textbf{176} (2012), no.~2, 865--923.
\url{https://doi.org/10.4007/annals.2012.176.2.4}.
Theorem~1.4, p.~872（Todd補正付きHKRとHochschild加群構造）。
\url{https://annals.math.princeton.edu/wp-content/uploads/annals-v176-n2-p04-p.pdf}.

\bibitem{HR}
D. Huybrechts and J. V. Rennemo,
\emph{Hochschild cohomology versus the Jacobian ring, and the Torelli theorem for cubic fourfolds},
Algebraic Geometry \textbf{6} (2019), 76--99.
\href{https://doi.org/10.14231/AG-2019-005}{doi:10.14231/AG-2019-005}.
Sections~2--3。一般のrotation付きHochschild環とJacobian環の比較の背景として参照する。
本稿の低次数の同定には次の\cite{LZ}を用いる。

\bibitem{LZ}
X. Lin and S. Zhang,
\emph{Serre algebra, matrix factorization and categorical Torelli theorem for hypersurfaces},
\href{https://arxiv.org/abs/2310.09927v2}{arXiv:2310.09927v2} (2024).
Theorem~2.2（dg比較とshift）、Theorem~3.5、Remark~3.6（sectorと積）、Section~5。
Proposition~3.7の互いに素という仮定はSerre代数への表示のための条件と区別する。

\bibitem{ExplicitUedaYoshinaga}
K. Ueda and M. Yoshinaga,
\emph{Logarithmic vector fields along smooth divisors in projective spaces},
Hokkaido Math. J. \textbf{38} (2009), 409--415.
\href{https://doi.org/10.14492/hokmj/1258553970}{doi:10.14492/hokmj/1258553970};
\href{https://arxiv.org/abs/0802.2133}{arXiv:0802.2133}.
Lemma~3：勾配イデアルとSebastiani--Thom型の関係。

\bibitem{ExplicitHwang}
J.-M. Hwang,
\emph{Symmetrizer group of a projective hypersurface},
J. Math. Soc. Japan \textbf{78} (2026), 55--62.
\href{https://doi.org/10.2969/jmsj/93469346}{doi:10.2969/jmsj/93469346};
\href{https://arxiv.org/abs/2501.08633}{arXiv:2501.08633}.
Theorems~1.5--1.6（勾配空間のfiberとsymmetrizer群）、
Proposition~2.2（可換代数の構造）。

\bibitem{PSPSmithTate}
S.~P.~Smith and J.~Tate,
\emph{The center of the 3-dimensional and 4-dimensional Sklyanin algebras},
K-Theory \textbf{8} (1994), 19--63, Section~1 (pp.~20--21).
\href{https://sites.math.washington.edu/~smith/Research/tate.pdf}{著者公開版}.
本稿では中心三次元による楕円曲線のtwisted homogeneous coordinate ringへの商を用いる。
有限位数を仮定する同論文後半の中心の分類は用いない。

\bibitem{PSPDeLaet}
K.~De Laet,
\emph{On the center of 3-dimensional and 4-dimensional Sklyanin algebras},
J. Algebra \textbf{487} (2017), 244--268,
\href{https://doi.org/10.1016/j.jalgebra.2017.05.032}{doi:10.1016/j.jalgebra.2017.05.032}.
参照箇所は\href{https://arxiv.org/abs/1612.06158}{arXiv:1612.06158v1}の
Definition~4.1、Sections~4.1--4.2、Theorem~4.4。
点曲線、平行移動$\tau=[a:b:c]$、中心三次元による点加群の消去を用いる。

\bibitem{TV} B. Toën and M. Vaquié,
\emph{Moduli of objects in dg-categories},
Ann. Sci. École Norm. Sup. (4) \textbf{40} (2007), 387--444.
\url{https://doi.org/10.1016/j.ansens.2007.05.001}.
Theorem 3.6, Corollaries 3.15, 3.17, 3.19, 3.21。
\url{https://www.math.univ-toulouse.fr/~vaquie/marque/vaquie_AnnEns40.pdf}.

\bibitem{ERLowen} W. Lowen,
\emph{Hochschild cohomology, the characteristic morphism and derived deformations},
Compos. Math. \textbf{144} (2008), 1557--1580.
\url{https://arxiv.org/abs/0707.2602v1}.
Theorem 4.8, Section 4.4。任意の小拡大に対するdg加群の記述は本文で説明する。

\bibitem{BFKER} M. Ballard, D. Favero and L. Katzarkov,
\emph{A category of kernels for equivariant factorizations and its implications for Hodge theory},
Publ. Math. Inst. Hautes Études Sci. \textbf{120} (2014), 1--111.
\url{https://doi.org/10.1007/s10240-013-0059-9}.
Theorems 5.39, 6.5。

\bibitem{HRER} D. Huybrechts and J. V. Rennemo,
\emph{Hochschild cohomology versus the Jacobian ring and the Torelli theorem for cubic fourfolds},
Algebraic Geometry \textbf{6} (2019), 76--99.
\url{https://doi.org/10.14231/AG-2019-005}.
Lemma 2.19, Sections 2.5, 3.3。

\bibitem{JLLZER} A. Jacovskis, X. Lin, Z. Liu and S. Zhang,
\emph{Infinitesimal categorical Torelli theorems for Fano threefolds},
J. Pure Appl. Algebra \textbf{227} (2023), 107418.
\url{https://arxiv.org/abs/2203.08187}.
Theorem 4.4(2)。

\bibitem{HuyER} D. Huybrechts,
\emph{The K3 category of a cubic fourfold},
Compos. Math. \textbf{153} (2017), 586--620.
\url{https://doi.org/10.1112/S0010437X16008137}.
Proposition 3.12, Section 3.5。

\bibitem{StacksER} The Stacks Project Authors, \emph{The Stacks Project}.
\url{https://stacks.math.columbia.edu/tag/02HX};
\url{https://stacks.math.columbia.edu/tag/055U}.
2026年11月13日確認。smoothnessの無限小判定とsmooth射のétale局所section。

\bibitem{ERLinZhang} X. Lin and S. Zhang,
\emph{Categorical Torelli and Fourier--Mukai partners of hypersurfaces},
arXiv:2611.05915v1 (2026).
\url{https://arxiv.org/abs/2611.05915}.
Section 2.1, Theorem 2.2, Proposition 2.4（元原稿の参照箇所）。
本稿の判定法の証明の入力ではなく、関連研究として挙げる。

\bibitem{AMSasha} Sasha,
Answer to \emph{Is the property of being a dg generator open?},
MathOverflow, 5 December 2016.
\url{https://mathoverflow.net/questions/256391/is-the-property-of-being-a-dg-generator-open}.
射影直線上の生成対象がskyscraperへ変形して生成性を失う例。

\bibitem{LP-AZ}
M. Artin and J. J. Zhang, \emph{Noncommutative projective schemes},
Adv. Math. \textbf{109} (1994), 228--287.
\url{https://doi.org/10.1006/aima.1994.1087}.

\bibitem{LP-BS}
J. Burke and G. Stevenson, \emph{The derived category of a graded Gorenstein ring},
arXiv:1507.00830v1 (2015).
\url{https://arxiv.org/abs/1507.00830}.
Section~4, Lemmas~6.1--6.2, Theorem~6.4およびその証明。

\bibitem{AMStacks}
The Stacks Project Authors, \emph{The Stacks Project}.
Tags \href{https://stacks.math.columbia.edu/tag/0BQT}{0BQT}
（有限線束生成子）および
\href{https://stacks.math.columbia.edu/tag/09M1}{09M1}
（compact対象とperfect対象）。

\bibitem{PHartshorne}
R.~Hartshorne,
\emph{Algebraic Geometry}, Graduate Texts in Mathematics 52,
Springer, 1977.
\href{https://doi.org/10.1007/978-1-4757-3849-0}{doi:10.1007/978-1-4757-3849-0}.

\bibitem{PEtingofLie}
P.~Etingof,
\emph{Lie Groups and Lie Algebras}, MIT lecture notes, 2024 course version.
\href{https://live.ocw.mit.edu/courses/18-755-lie-groups-and-lie-algebras-ii-spring-2024/mit18_755_s24_lec_full.pdf}{MIT OpenCourseWare 公開原稿}.
本文では Proposition~16.9、\S19.4、Theorem~25.17、\S32.2 を参照する。

\bibitem{PMoriSmith}
I.~Mori and S.~P.~Smith,
The classification of 3-Calabi--Yau algebras with 3 generators and 3 quadratic relations,
\href{https://arxiv.org/abs/1502.07403v2}{arXiv:1502.07403v2} (2016).
Theorems~1.3, 1.4, 4.2 and Corollaries~1.6, 2.7, 2.9.
点曲線のHessian公式はLemma~7.1, Theorem~7.2, pp.~23--25。

\bibitem{PBergerTaillefer}
R.~Berger and R.~Taillefer,
Poincar\'e--Birkhoff--Witt deformations of Calabi--Yau algebras,
\emph{Journal of Noncommutative Geometry} \textbf{1} (2007), 241--270.
\href{https://arxiv.org/abs/math/0610112}{arXiv:math/0610112}.
本文で用いる potential による PBW 変形は Theorem~1.1(i)。

\bibitem{PKontsevich}
M.~Kontsevich,
Deformation quantization of Poisson manifolds,
\emph{Letters in Mathematical Physics} \textbf{66} (2003), 157--216.
\href{https://doi.org/10.1023/B:MATH.0000027508.00421.bf}{doi:10.1023/B:MATH.0000027508.00421.bf};
原稿名は \emph{Deformation quantization of Poisson manifolds, I},
\href{https://arxiv.org/abs/q-alg/9709040}{arXiv:q-alg/9709040} (1997).
参照箇所は原稿のSections~4.4--4.6（Maurer--Cartan変形functorとformality）、
6.4（$L_\infty$関係）、8.2および8.3.3.1（中心のtangent mapとwheel）。

\bibitem{PDolgushev}
V.~Dolgushev,
The Van den Bergh duality and the modular symmetry of a Poisson variety,
\href{https://arxiv.org/abs/math/0612288v3}{arXiv:math/0612288v3} (2007).
Theorem~3 および Proposition~1。

\bibitem{BondalGualtieriPym}
M.~Gualtieri and B.~Pym,
\emph{Poisson modules and degeneracy loci},
Proc. Lond. Math. Soc. \textbf{107} (2013), 627--654.
\href{https://doi.org/10.1112/plms/pds090}{doi:10.1112/plms/pds090};
\href{https://arxiv.org/abs/1203.4293}{arXiv:1203.4293}.
Bondal予想の定式化、四次元Fanoの場合の証明。

\bibitem{BondalDruelPereiraPymTouzet}
S.~Druel, J.~V.~Pereira, B.~Pym and F.~Touzet,
\emph{Bondal's conjecture in dimension five},
\href{https://arxiv.org/abs/2606.03424}{arXiv:2606.03424} (2026).
Fano五次元の全場合、一般の奇数次元の部分結果。

\bibitem{FuHwangCI}
B. Fu and J.-M. Hwang,
\emph{Isotrivial VMRT-structures of complete intersection type},
Asian Journal of Mathematics \textbf{22} (2018), 333--354.
Lemma 4.1 (Br\"uckmannの消滅定理を引用).
\url{https://doi.org/10.4310/AJM.2018.v22.n2.a10}.

\bibitem{BorelRemmert}
A. Borel and R. Remmert,
\emph{\"Uber kompakte homogene K\"ahlersche Mannigfaltigkeiten},
Mathematische Annalen \textbf{145} (1962), 429--439, Theorem I.
\url{https://wglobke.github.io/publications/BorelRemmert1962.pdf}.

\bibitem{EG}
P. Etingof and V. Ginzburg,
\emph{Noncommutative del Pezzo surfaces and Calabi--Yau algebras},
Journal of the European Mathematical Society \textbf{12} (2010), 1371--1416.
Sections 3.5--3.6, Theorem 3.6.2, Proposition 8.1.1(i) およびその証明。
\url{https://doi.org/10.4171/JEMS/235}.
\href{https://arxiv.org/abs/0709.3593v4}{arXiv:0709.3593v4}.

\bibitem{LBSVNE12} L. Le Bruyn, S. P. Smith and M. Van den Bergh,
\emph{Central extensions of three dimensional Artin--Schelter regular algebras},
Math. Z. \textbf{222} (1996), 171--212.
Theorem 2.6, Corollary 2.7, Remark 2.8, Theorem 3.1.3.
\url{https://sites.math.washington.edu/~smith/Research/Central-extns.pdf}.

\bibitem{MSNE12} I. Mori and S. P. Smith,
\emph{The classification of 3-Calabi--Yau algebras with 3 generators and 3 quadratic relations},
Math. Z. \textbf{287} (2017), 215--241.
Lemma 4.1, Theorem 4.2, Corollary 4.3; arXiv v2, pp. 15--17.
\url{https://arxiv.org/abs/1502.07403}; doi:10.1007/s00209-016-1824-5.

\bibitem{BergNE12} G. M. Bergman,
\emph{The diamond lemma for ring theory}, Adv. Math. \textbf{29} (1978), 178--218, Theorem 1.2.
\url{https://doi.org/10.1016/0001-8708(78)90010-5}.

\bibitem{BTNE12} R. Berger and R. Taillefer,
\emph{Poincar\'e--Birkhoff--Witt deformations of Calabi--Yau algebras},
J. Noncommut. Geom. \textbf{1} (2007), 241--270.
Theorems 3.1 and 3.6, especially the potential bimodule complex in the proof of Theorem 3.6.
\url{https://arxiv.org/abs/math/0610112}; doi:10.4171/JNCG/6.

\bibitem{SmithTate}
S. P. Smith and J. Tate,
\emph{The center of the 3-dimensional and 4-dimensional Sklyanin algebras},
K-Theory \textbf{8} (1994), 19--63.
Corollary 2.3 (p.~26); Theorem 3.7(c) (p.~38).
\url{https://doi.org/10.1007/BF00962089};
\url{https://sites.math.washington.edu/~smith/Research/tate.pdf}.

\bibitem{DeLaetLeBruyn}
K. De Laet and L. Le Bruyn,
\emph{The geometry of representations of 3-dimensional Sklyanin algebras},
Algebras and Representation Theory \textbf{18} (2015), 761--776.
arXiv:1405.1158v1, Theorem 1 and Section 2 (canonical central cubic formula);
the result is a reformulation of the Artin--Schelter and Smith--Tate formula.
\url{https://arxiv.org/abs/1405.1158}.

\bibitem{ArtebaniDolgachev}
M. Artebani and I. Dolgachev,
\emph{The Hesse pencil of plane cubic curves},
L'Enseignement Math\'ematique (2) \textbf{55} (2009), no.~3--4, 235--273,
Lemma 1 (every smooth plane cubic admits a Hesse canonical form).
\url{https://doi.org/10.4171/LEM/55-3-3};
\url{https://arxiv.org/abs/math/0611590}.

\bibitem{LPTPoissonCubic}
F.~Loray, J.~V.~Pereira and F.~Touzet,
\emph{Foliations with trivial canonical bundle on Fano 3-folds},
Math. Nachr. \textbf{286} (2013), 921--940.
\href{https://arxiv.org/abs/1112.3871}{arXiv:1112.3871}.

\bibitem{PStacksQuotient}
The Stacks Project Authors,
\emph{The Stacks Project},
Criteria for Representability, Theorem~97.17.2
(Tag~\texttt{06FI}), flat locally finitely presented groupoids and algebraic quotient stacks.
\url{https://stacks.math.columbia.edu/tag/06FI}.

\bibitem{PGLefschetz}
A.~Grothendieck,
\emph{Cohomologie locale des faisceaux coh\'erents et th\'eor\`emes
de Lefschetz locaux et globaux (SGA~2)},
Expos\'e~XII, Corollaire~3.7 (Picard groups of sufficiently
high-dimensional smooth projective complete intersections).
\url{https://grothendiecksga.com/read/sga2/en/XII-3.html}.

\bibitem{PMatsumuraMonsky}
H.~Matsumura and P.~Monsky,
\emph{On the automorphisms of hypersurfaces},
J. Math. Kyoto Univ. \textbf{3} (1963/64), 347--361.
\url{https://doi.org/10.1215/kjm/1250524785}.

\bibitem{PAZAmbient}
M.~Artin and J.~J.~Zhang,
\emph{Noncommutative projective schemes},
Adv. Math. \textbf{109} (1994), 228--287,
\href{https://doi.org/10.1006/aima.1994.1087}{doi:10.1006/aima.1994.1087}.

\bibitem{LP-Keller} B. Keller, \emph{On differential graded categories},
International Congress of Mathematicians, Vol. II (2006), 151--190, Theorem 3.8.
\url{https://arxiv.org/abs/math/0601185}.

\bibitem{LP-Orlov} D. Orlov,
\emph{Smooth and proper noncommutative schemes and gluing of DG categories},
Adv. Math. \textbf{302} (2016), 59--105, Section 3.4, Proposition 3.31.
\url{https://arxiv.org/abs/1402.7364}.

\bibitem{PetersenTransfer} D.~Petersen,
\emph{A closer look at Kadeishvili's theorem},
Higher Structures \textbf{4} (2020), no.~2, 211--221.
\url{https://doi.org/10.21136/HS.2020.14}.

\bibitem{BurkeTransfer} J.~Burke,
\emph{Transfer of $A_\infty$ structures to projective resolutions},
arXiv:1801.08933 (2018).
\url{https://arxiv.org/abs/1801.08933}.

\bibitem{AinftyPerry2019} A.~Perry,
\emph{Noncommutative homological projective duality},
Advances in Mathematics \textbf{350} (2019), 877--972,
Lemma 4.15(3).
\url{https://doi.org/10.1016/j.aim.2019.04.052}.

\bibitem{UeyamaIsolated}
K. Ueyama,
\emph{Noncommutative graded Gorenstein isolated singularities},
in \emph{Proceedings of the 45th Symposium on Ring Theory and Representation Theory}
(Shinshu University, September 7--9, 2012),
K. Koike (ed.), February 2013, 158--163.
Theorem 7, printed p.~161.
\url{https://ring-theory-japan.com/ring/oldmeeting/2012/report2012/proceedings.pdf}.

\bibitem{ZhangTwist}
J.~J. Zhang,
\emph{Twisted graded algebras and equivalences of graded categories},
Proceedings of the London Mathematical Society (3) \textbf{72} (1996), no.~2, 281--311.
Theorems 1.1 and 1.3.
\url{https://doi.org/10.1112/plms/s3-72.2.281}.
\url{https://deepblue.lib.umich.edu/bitstream/handle/2027.42/135651/plms0281.pdf}.

\bibitem{RRZ}
M. Reyes, D. Rogalski and J.~J. Zhang,
\emph{Skew Calabi--Yau algebras and homological identities},
Advances in Mathematics \textbf{264} (2014), 308--354.
Theorems 0.3 and 5.4; Example 5.5.
\url{https://doi.org/10.1016/j.aim.2014.07.010}.
\url{https://arxiv.org/abs/1302.0437}.

\bibitem{Griffiths}
P.~A. Griffiths,
\emph{On the periods of certain rational integrals.~I},
Annals of Mathematics (2) \textbf{90} (1969), no.~3, 460--495.
Section 8, especially Theorems 8.1 and 8.3 and formula (8.6).
\url{https://doi.org/10.2307/1970746}.
\url{https://www.cimat.mx/~luis/seminarios/Teoria-de-Hodge/period_integral_I.pdf}.

\bibitem{KuznetsovFCY}
A. Kuznetsov,
\emph{Calabi--Yau and fractional Calabi--Yau categories},
Journal f\"ur die reine und angewandte Mathematik \textbf{753} (2019), 239--267.
Lemmas 2.6, 3.13 and 3.19; Corollary 4.1.
\url{https://doi.org/10.1515/crelle-2017-0004}.

\bibitem{CHL}
C.-H. Cho, H. Hong and S.-C. Lau,
\emph{Noncommutative homological mirror functor},
Memoirs of the American Mathematical Society \textbf{271} (2021), no.~1326.
Chapter 5, especially Theorems 5.15 and 5.16.
\url{https://doi.org/10.1090/memo/1326}.
\url{https://arxiv.org/abs/1512.07128}.

\bibitem{Lee}
S. Lee,
\emph{Noncommutative homological mirror symmetry of elliptic curves},
Kyoto Journal of Mathematics \textbf{61} (2021), no.~3, 723--743.
Theorem 6.3 and the corollary in the introduction.
\url{https://arxiv.org/abs/2004.09145}.

\bibitem{CisinskiTabuada}
D.-C. Cisinski and G. Tabuada,
\emph{Lefschetz and Hirzebruch--Riemann--Roch formulas via noncommutative motives},
Journal of Noncommutative Geometry \textbf{8} (2014), 1171--1190.
Section 4 and Example 5.6.
\url{https://doi.org/10.4171/JNCG/183}.
\url{https://www.math.univ-toulouse.fr/~dcisinsk/ncLHRR.pdf}.

\bibitem{MathewHH}
A. Mathew,
\emph{Kaledin's degeneration theorem and topological Hochschild homology},
Geometry \& Topology \textbf{24} (2020), no.~6, 2675--2708.
Theorem 1.3.
\url{https://doi.org/10.2140/gt.2020.24.2675}.
\url{https://arxiv.org/abs/1710.09045}.

\bibitem{PerryIHC}
A. Perry,
\emph{The integral Hodge conjecture for two-dimensional Calabi--Yau categories},
Compositio Mathematica \textbf{158} (2022), no.~2, 287--333.
Lemmas 3.4 and 3.6; Theorem 3.5.
\url{https://doi.org/10.1112/S0010437X22007266}.

\bibitem{Beilinson}
A. A. Beilinson,
\emph{Coherent sheaves on \(\PP^n\) and problems of linear algebra},
Functional Analysis and Its Applications \textbf{12} (1978), 214--216.
\url{https://doi.org/10.1007/BF01681436}.
原著Section 4の対角Koszul分解を確認。
\url{https://www.mathnet.ru/eng/faa2009}.

\bibitem{DDL}
O. De Deken and W. Lowen,
\emph{Abelian and derived deformations in the presence of \(\ZZ\)-generating geometric helices},
Journal of Noncommutative Geometry \textbf{5} (2011), 477--505.
Theorem 4.5; Propositions 5.5, 5.8, 5.9; Lemma 5.14.
\url{https://doi.org/10.4171/JNCG/83}.

\bibitem{Keller}
B. Keller, \emph{Deriving DG categories},
Annales scientifiques de l'École Normale Supérieure \textbf{27} (1994), 63--102.
\url{https://doi.org/10.24033/asens.1689}.

\bibitem{Mathew}
A. Mathew, \emph{The Galois group of a stable homotopy theory},
Advances in Mathematics \textbf{291} (2016), 403--541.
Proposition 3.28, Remark 3.29, Proposition 3.35.
\url{https://arxiv.org/abs/1404.2156}.

\bibitem{CCKM}
T. Cassidy, A. Conner, E. Kirkman and W. F. Moore,
\emph{Periodic free resolutions from twisted matrix factorizations},
Journal of Algebra \textbf{455} (2016), 137--163.
Proposition 5.7, Theorem 5.8.
\url{https://doi.org/10.1016/j.jalgebra.2016.01.037}.

\bibitem{BFK}
M. Ballard, D. Favero and L. Katzarkov,
\emph{A category of kernels for equivariant factorizations and its implications for Hodge theory},
Publications Mathématiques de l'IHÉS \textbf{120} (2014), 1--111.
Theorem 5.39のKoszul cohomologyによる公式、Theorem 5.15とCorollary 5.18のequivariant外部積.
\url{https://doi.org/10.1007/s10240-013-0059-9}.
確認した公刊本文 \url{https://www.numdam.org/article/PMIHES_2014__120__1_0.pdf}.

\bibitem{Huy}
D. Huybrechts, \emph{The K3 category of a cubic fourfold},
Compositio Mathematica \textbf{153} (2017), 586--620.
Proposition 3.12およびSection 3.5.
\url{https://doi.org/10.1112/S0010437X16008137}.

\bibitem{Lurie}
J. Lurie, \emph{Derived Algebraic Geometry X: Formal Moduli Problems},
2011. Theorems 5.1.9, 5.3.16およびProposition 5.3.21.
\url{https://www.math.ias.edu/~lurie/papers/DAG-X.pdf}.

\bibitem{BKP}
A. Blanc, L. Katzarkov and P. Pandit,
\emph{Generators in formal deformations of categories},
Compositio Mathematica \textbf{154} (2018).
本文で用いるのはarXiv版のDefinition 3.3およびSection 3.3.
\url{https://arxiv.org/abs/1705.00655}.

\bibitem{KL}
B. Keller and W. Lowen,
\emph{On Hochschild cohomology and Morita deformations},
International Mathematics Research Notices (2009).
\url{https://webusers.imj-prg.fr/~bernhard.keller/publ/KellerLowenHochDefo.pdf}.

\bibitem{Walton}
C. Walton, \emph{Representation theory of three-dimensional Sklyanin algebras},
Nuclear Physics B \textbf{860} (2012), 167--185.
Theorem 2.14および式(2.15)--(2.16).
\url{https://doi.org/10.1016/j.nuclphysb.2012.02.015}.

\bibitem{CGWZ}
K. Chan, J. Gaddis, R. Won and J. J. Zhang,
\emph{Ozone groups of Artin--Schelter regular algebras satisfying a polynomial identity},
arXiv:2312.13014, v2 (2024). Introductionの非退化Sklyanin代数の記述。
\url{https://arxiv.org/abs/2312.13014}.

\bibitem{Elagin}
A. Elagin, \emph{On equivariant triangulated categories},
arXiv:1403.7027, Theorem 4.2.
\url{https://arxiv.org/abs/1403.7027}.

\bibitem{ELS}
A. Elagin, V. A. Lunts and O. M. Schn\"urer,
\emph{Smoothness of Derived Categories of Algebras},
Moscow Mathematical Journal \textbf{20} (2020), no. 2, 277--309.
確認した一次資料はarXiv:1810.07626v2 (22 March 2019),
Theorem 4.4 (Verdier局所化), Proposition 4.10とTheorem 4.15 (稠密開集合と生成子),
Theorem 4.18 (外部積生成), Theorem 5.1 (有限中心代数の有界導来圏のsmoothness).
\url{https://arxiv.org/abs/1810.07626};
\url{https://www.mathjournals.org/mmj/2020-020-002/2020-020-002-003.html}.

\bibitem{Stacks}
The Stacks Project Authors, \emph{The Stacks Project}.
Tag 02HX (Lemma 37.12.1, Noetherian底上の無限小smoothness判定),
Tag 055U (Lemma 37.38.6, smooth morphismのétale局所section).
\url{https://stacks.math.columbia.edu/tag/02HX};
\url{https://stacks.math.columbia.edu/tag/055U}.
参照日2026年11月4日。

\bibitem{Pirozhkov}
D. Pirozhkov, \emph{A categorical Torelli theorem for hypersurfaces},
Bulletin of the London Mathematical Society \textbf{56} (2024), 3075--3089.
\url{https://doi.org/10.1112/blms.13117}.
公刊版のabstractで、cubicに対する仮定 \(n>3\) を確認。
本稿の全次元の一意性には \cite[Theorem 1.7]{LZ} を用いる。

\bibitem{HuyUpdate}
D. Huybrechts, \emph{The K3 category of a cubic fourfold: an update},
Beiträge zur Algebra und Geometrie / Contributions to Algebra and Geometry (2025).
Section 9: Unknown deformations.
\url{https://doi.org/10.1007/s13366-025-00788-9}.

\bibitem{ChanMulti}
D. Chan, \emph{Twisted multi-homogeneous coordinate rings},
Journal of Algebra \textbf{223} (2000), 438--456.
Theorem 3.4; proof of Theorem 5.2, Claim 5.3 and Lemmas 5.4--5.6;
Corollaries 5.8--5.9.
\url{https://doi.org/10.1006/jabr.1999.7938}.
確認した著者版 \url{https://web.maths.unsw.edu.au/~danielch/paper/tbcr.ps}, pp.8--9,14--17.

\bibitem{KeelerMulti}
D. S. Keeler, \emph{Noncommutative ampleness for multiple divisors},
Journal of Algebra \textbf{265} (2003), no. 1, 299--311.
Theorem 2.7, Corollary 2.8 (任意有限個のbimoduleのNC-ampleness);
Hypothesis 3.1, Theorem 3.3, Corollary 3.5 (二因子のnoetherian性).
番号はarXiv一次TeXソースで確認し、刊行版との番号一致は未照合。
\url{https://doi.org/10.1016/S0021-8693(03)00126-1};
\url{https://arxiv.org/abs/math/0210417};
\url{https://arxiv.org/e-print/math/0210417}.

\bibitem{WWY3}
C. Walton, X. Wang and M. Yakimov,
\emph{Poisson geometry of PI three-dimensional Sklyanin algebras},
Proceedings of the London Mathematical Society \textbf{118} (2019), 1471--1500.
Theorem 1.3(1),(2a), verified in arXiv:1704.04975v3, p.~3.
\url{https://doi.org/10.1112/plms.12220};
\url{https://arxiv.org/abs/1704.04975}.

\bibitem{ToenFinitude}
B. Toën, \emph{Finitude homotopique des dg-algèbres propres et lisses},
Proceedings of the London Mathematical Society (3) \textbf{98} (2009), 217--240.
Lemmas 3.13--3.14. Lemma 3.14の結論は部分前層のquasi-compactnessであり、一般の開性ではない。
\url{https://doi.org/10.1112/plms/pdn027};
確認した著者版 \url{https://www.math-beb.fr/wp-content/uploads/2026/10/finitesatdgalg.pdf}.

\bibitem{SnowdenDVR}
A. Snowden, \emph{Lecture 8: Elliptic curves over DVRs},
Math 679 lecture notes, University of Michigan, 2013.
Section ``Reduction of torsion points'', good reductionでのprime-to-residue-characteristic torsionの還元の単射性。
\url{https://websites.umich.edu/~asnowden/teaching/2013/679/L08.html}.
参照日2026年11月6日。

\bibitem{HMY}
J.-W. He, X.-C. Ma and Y. Ye, \emph{Generalized Kn\"orrer's Periodicity Theorem},
arXiv:2107.06438v2 (16 August 2021).
対象は中心二次元の商であり、本稿の三次の場合には直接適用しない。
\url{https://arxiv.org/abs/2107.06438}.

\bibitem{OrlovGlue}
D. Orlov,
\emph{Smooth and proper noncommutative schemes and gluing of DG categories},
Advances in Mathematics \textbf{302} (2016), 59--105.
Theorem 3.25 (smoothness of dg gluing; Proposition 3.24はsmoothnessからregularityのみ)。
\url{https://doi.org/10.1016/j.aim.2016.07.014};
\url{https://arxiv.org/abs/1402.7364}.

\bibitem{BurkeStevenson}
J. Burke and G. Stevenson,
\emph{The derived category of a graded Gorenstein ring},
arXiv:1507.00830v1 (2015), Definition 5.1 and Theorem 6.4.
\url{https://arxiv.org/abs/1507.00830}.

\bibitem{MBNegron}
C. Negron, \emph{The cup product on Hochschild cohomology via twisting cochains and applications to Koszul rings},
\href{https://arxiv.org/abs/1304.0527v4}{arXiv:1304.0527v4} (2016), Theorem~5.3.

\bibitem{MBTu}
J. Tu, \emph{Matrix factorizations via Koszul duality},
Compos. Math. \textbf{150} (2014), 1549--1578;
\href{https://arxiv.org/abs/1009.4151v2}{arXiv:1009.4151v2}, Appendix~A, Lemma~A.1.

\bibitem{MBLowen}
W. Lowen, \emph{Hochschild cohomology, the characteristic morphism and derived deformations},
Compos. Math. \textbf{144} (2008), 1557--1580;
\href{https://arxiv.org/abs/0707.2602}{arXiv:0707.2602}, Propositions~3.11, 3.13.

\bibitem{MBGHLR}
A. Guan, J. Holstein, A. Lazarev and K. Rasmussen,
\emph{Hochschild cohomology of the second kind: Koszul duality and Morita invariance},
\href{https://arxiv.org/abs/2312.16645v5}{arXiv:2312.16645v5} (2026),
Theorems~4.7, 6.14 and Remark~6.16.

\bibitem{MBLinZhang}
X. Lin and S. Zhang, \emph{Serre algebra, matrix factorization and categorical Torelli theorem for hypersurfaces},
\href{https://arxiv.org/abs/2310.09927v2}{arXiv:2310.09927v2} (2024),
Theorem~3.5 and Remark~3.6. 本稿では体上の特殊ファイバーの同定に用いる。

\bibitem{MBStacks}
The Stacks Project Authors, \emph{The Stacks Project},
\href{https://stacks.math.columbia.edu/tag/00T7}{Tag 00T7}.
\'Etale近傍の構成には相対Jacobianの可逆性を用いる。

\bibitem{FARStacks}
The Stacks Project Authors,
\emph{The Stacks Project}, Tags 00MC and 0BNH,
\url{https://stacks.math.columbia.edu/tag/00MC},
\url{https://stacks.math.columbia.edu/tag/0BNH}.
完備化の忠実平坦性、有限加群の完備化。2026年11月10日確認。

\bibitem{FARLinZhang}
X. Lin and S. Zhang,
\emph{Serre algebra, matrix factorization and categorical Torelli theorem for hypersurfaces},
arXiv:2310.09927v2 (2024), Theorem 3.5, Remark 3.6, Section 5.
\url{https://arxiv.org/html/2310.09927v2}.

\bibitem{FARHwang}
J.-M. Hwang,
\emph{Symmetrizer group of a projective hypersurface},
J. Math. Soc. Japan \textbf{78} (2026), 55--62;
arXiv:2501.08633v1, Sections 1--3.
\url{https://arxiv.org/html/2501.08633v1}.

\bibitem{GCDBFK}
M.~Ballard, D.~Favero and L.~Katzarkov,
\emph{Orlov spectra: bounds and gaps},
arXiv:1012.0864v3 (2013), Section~2 and Theorem~5.15.
\url{https://arxiv.org/abs/1012.0864}.
参照したPDFの第41頁の定理と第6頁の基礎体・生成段数の規約を確認。

\bibitem{CGTu}
J.~Tu,
\emph{Matrix factorizations via Koszul duality},
Compositio Mathematica \textbf{150} (2014), 1549--1578.
\url{https://doi.org/10.1112/S0010437X14007295}.
\url{https://arxiv.org/abs/1009.4151}.
参照箇所はarXiv v2のSections~2.2, 6.1およびTheorem~6.1。
curved twisted complexとgradingの枠組みとして参照し、
非可換Artin基底上の本稿の比較定理全体を引用するものではない。

\bibitem{ROGStevenson}
G.~Stevenson, \emph{Subcategories of singularity categories via tensor actions},
\href{https://arxiv.org/abs/1105.4698v3}{arXiv:1105.4698v3},
Theorem~6.13 (PDF p.~26),
\href{https://doi.org/10.1112/S0010437X1300746X}{doi:10.1112/S0010437X1300746X}.
原論文の非次数付き分類を引用し、次数シフト軌道への移行は本文で証明する。

\bibitem{ROGDyckerhoff}
T.~Dyckerhoff, \emph{Compact generators in categories of matrix factorizations},
Duke Math. J. \textbf{159} (2011), 223--274;
\href{https://arxiv.org/abs/0904.4713v5}{arXiv:0904.4713v5},
\S2.4, \S5.6, Theorem~5.9. 本文の三次Euler計算は直接も証明する。

\bibitem{ROGStacks}
The Stacks Project Authors, \emph{The Stacks Project},
\href{https://stacks.math.columbia.edu/tag/09SM}{Tag 09SM},
\S13.37, Lemmas~13.37.3--4 and Proposition~13.37.6;
\href{https://stacks.math.columbia.edu/tag/00FV}{Tag 00FV}, Theorem~10.34.1;
\href{https://stacks.math.columbia.edu/tag/0BQQ}{Tag 0BQQ}, \S36.16 (projective subscheme generators).

\bibitem{ROGTV}
B.~To\"en and M.~Vaqui\'e, \emph{Moduli of objects in dg-categories},
\href{https://arxiv.org/abs/math/0503269v5}{arXiv:math/0503269v5},
Lemmas~2.8 and 2.10 (PDF pp.~21--24).
片側perfect性の比較と、perfect対象・射のfiltered-colimit降下に用いる。

\bibitem{RASBFK}
M. Ballard, D. Favero and L. Katzarkov,
\emph{A category of kernels for equivariant factorizations and its implications for Hodge theory},
Publ. Math. IH\'ES \textbf{120} (2014), 1--111.
\href{https://doi.org/10.1007/s10240-013-0059-9}{doi:10.1007/s10240-013-0059-9};
\href{https://arxiv.org/abs/1105.3177}{arXiv:1105.3177}.

\bibitem{RASDyck}
T. Dyckerhoff,
\emph{Compact generators in categories of matrix factorizations},
Duke Math. J. \textbf{159} (2011), 223--274.
\href{https://arxiv.org/abs/0904.4713}{arXiv:0904.4713}.

\bibitem{RASLZ}
X. Lin and S. Zhang,
\emph{Serre algebra, matrix factorization and categorical Torelli theorem for hypersurfaces},
\href{https://arxiv.org/html/2310.09927v2}{arXiv:2310.09927v2},
Theorem~3.5 and Remark~3.6.

\bibitem{KJDito}
G. Dito,
\emph{Kontsevich star-product on the dual of a Lie algebra},
Letters in Mathematical Physics \textbf{48} (1999), 307--322.
\href{https://doi.org/10.1023/A:1007643618406}{doi:10.1023/A:1007643618406};
\href{https://arxiv.org/abs/math/9905080}{arXiv:math/9905080}.
式~(3)（arXiv版の印刷頁5）は一般Poisson tensorに対する二次星積とeye項の係数を与える。
本稿では同論文のbivectorを$\pi/2$に置き換える。Section~3はWeyl対称化を扱う。

\bibitem{KJShoikhet}
B. Shoikhet,
\emph{Vanishing of the Kontsevich integrals of the wheels},
\href{https://arxiv.org/abs/math/0007080}{arXiv:math/0007080} (2000).
Section~1.3はboundary vertexのないtangent mapの偶数wheelの重み$W_{2k}=0$を述べる。
本稿では$W_2=0$を用い、星積の二つのboundary vertexを持つeye項とは区別する。

\bibitem{MNFDolgushev}
V.~Dolgushev,
\emph{Covariant and equivariant formality theorems},
Advances in Mathematics \textbf{191} (2005), 147--177.
\href{https://doi.org/10.1016/j.aim.2004.02.001}{doi:10.1016/j.aim.2004.02.001};
\href{https://arxiv.org/abs/math/0307212}{arXiv:math/0307212}.
有限群不変な接続を使うformalityの同変化。本文のaffine線形作用ではKontsevichのGL同変性を使う。

\bibitem{AROBFK}
M.~Ballard, D.~Favero and L.~Katzarkov,
\emph{A category of kernels for equivariant factorizations and its implications for Hodge theory},
Publ. Math. Inst. Hautes Études Sci. \textbf{120} (2014), 1--111,
\href{https://doi.org/10.1007/s10240-013-0059-9}{doi:10.1007/s10240-013-0059-9},
Theorem~5.39 (printed pp.~82--84).

\bibitem{AROLZ}
X.~Lin and S.~Zhang,
\emph{Serre algebra, matrix factorization and categorical Torelli theorem for hypersurfaces},
\href{https://arxiv.org/abs/2310.09927v2}{arXiv:2310.09927v2},
Theorem~3.5, Remark~3.6.

\bibitem{CDGHLR}
A.~Guan, J.~Holstein, A.~Lazarev and K.~Rasmussen,
\emph{Hochschild cohomology of the second kind: Koszul duality and Morita invariance},
arXiv:2312.16645v5 (2026-09-17).
\href{https://arxiv.org/abs/2312.16645v5}{arXiv:2312.16645v5}.
Theorem~4.7, Theorem~6.14 and Remark~6.16; pp.~32--33 of the PDF.
本稿では包絡代数上の導来圏と曲率付き双加群の導来圏の区別に用いる。

\bibitem{OSTV}
B.~To\"en and M.~Vaqui\'e,
\emph{Moduli of objects in dg-categories},
Ann. Sci. \`Ecole Norm. Sup. (4) \textbf{40} (2007), 387--444.
\href{https://doi.org/10.1016/j.ansens.2007.05.001}{doi:10.1016/j.ansens.2007.05.001};
\href{https://arxiv.org/abs/math/0503269}{arXiv:math/0503269}.
Lemma~2.8(2),(3), Lemma~2.10 (printed pp.~404--407).

\bibitem{OSPopescu}
The Stacks Project Authors, \emph{The Stacks Project},
\href{https://stacks.math.columbia.edu/tag/07GC}{Tag 07GC},
Theorem~16.12.1: regular ring maps are filtered colimits of smooth maps.

\bibitem{OSEtale}
The Stacks Project Authors, \emph{The Stacks Project},
\href{https://stacks.math.columbia.edu/tag/07M7}{Tag 07M7}, Lemma~15.9.14:
smooth代数の剰余切断の\'etale持上げ。

\bibitem{OSStacksSplit}
The Stacks Project Authors, \emph{The Stacks Project},
\href{https://stacks.math.columbia.edu/tag/068U}{Tag 068U}, Lemma~15.78.4:
特殊ファイバーのcohomology消滅によるperfect複体の局所分裂。

\bibitem{OSApprox}
The Stacks Project Authors, \emph{The Stacks Project},
\href{https://stacks.math.columbia.edu/tag/07QZ}{Tag 07QZ}, Theorem~16.13.2:
G-ring上の形式解の\'etale近似。本文はそのPopescuと切断による証明法を用いる。

\bibitem{CRSBS}
J.~Burke and G.~Stevenson, \emph{The derived category of a graded Gorenstein ring},
\href{https://arxiv.org/abs/1507.00830}{arXiv:1507.00830v1} (2015),
Proposition~4.18, Lemmas~6.1--6.2 and Theorem~6.4.
PDF pp.~12--14, 17--19を確認。

\bibitem{CRSTV}
B.~To\"en and M.~Vaqui\'e, \emph{Moduli of objects in dg-categories},
Ann. Sci. \'Ecole Norm. Sup. (4) \textbf{40} (2007), 387--444;
\href{https://arxiv.org/abs/math/0503269}{arXiv:math/0503269v5}, Lemma~2.8(2).
PDF pp.~21--22を確認。pseudo-perfectからperfectへの原理を用いる。

\bibitem{CRSPerry}
A.~Perry, \emph{Noncommutative homological projective duality},
Adv. Math. \textbf{350} (2019), 877--972, Lemma~4.15(3).
\href{https://www.sciencedirect.com/science/article/am/pii/S000187081930221X}{著者受理稿}。
admissible半直交成分とsmoothnessの同値を用いる。

\bibitem{CRSStevenson}
G.~Stevenson, \emph{Subcategories of singularity categories via tensor actions},
Compos. Math. \textbf{150} (2014), 229--272;
\href{https://arxiv.org/abs/1105.4698}{arXiv:1105.4698v3}, Theorem~6.13 (PDF p.~26).
今回再確認したのはaffine hypersurfaceの全thick部分圏の分類である。

\bibitem{DKITu}
J. Tu, \emph{Matrix factorizations via Koszul duality},
Compos. Math. \textbf{150} (2014), 1549--1578;
\href{https://arxiv.org/abs/1009.4151}{arXiv:1009.4151}, Appendix~A, Lemma~A.1.

\bibitem{DKIBFK}
M. Ballard, D. Favero and L. Katzarkov,
\emph{A category of kernels for equivariant factorizations and its implications for Hodge theory},
Publ. Math. Inst. Hautes \`Etudes Sci. \textbf{120} (2014), 1--111;
\href{https://doi.org/10.1007/s10240-013-0059-9}{doi:10.1007/s10240-013-0059-9},
Theorem~5.22, Remark~5.21.

\bibitem{DKIDyckerhoff}
T. Dyckerhoff, \emph{Compact generators in categories of matrix factorizations},
Duke Math. J. \textbf{159} (2011), 223--274;
\href{https://arxiv.org/abs/0904.4713}{arXiv:0904.4713}, Section~5.

\bibitem{TWCrainic}
M. Crainic,
\emph{On the perturbation lemma, and deformations},
arXiv:math/0403266.
\url{https://arxiv.org/abs/math/0403266}.
著者公開版\url{https://www.staff.science.uu.nl/~crain101/pert2.pdf}の
Theorem 2.3とRemark 2.4（2026年11月12日本文確認）。
同版の$ip=1+bh+hb$という符号は、本稿の$h=-s$に対応する。

\bibitem{RDPP}
A.~Polishchuk and L.~Positselski,
\emph{Hochschild (co)homology of the second kind I},
Trans. Amer. Math. Soc. \textbf{364} (2012), no.~10, 5311--5368.
\href{https://doi.org/10.1090/S0002-9947-2012-05667-4}{doi:10.1090/S0002-9947-2012-05667-4};
\href{https://arxiv.org/abs/1010.0982v2}{arXiv:1010.0982v2}.
\S3.3では有限射影次元を持つ対象の第二種Extを扱う。
\S3.5の対角比較条件を、本稿の非可換族について確認済みとはしない。

\bibitem{GRToen}
B.~To\"en, \emph{The homotopy theory of dg-categories and derived Morita theory},
Invent. Math. \textbf{167} (2007), 615--667.
\href{https://arxiv.org/abs/math/0408337}{arXiv:math/0408337};
\S2 (arbitrary commutative base ring), Theorems~8.9 and 8.15.

\bibitem{GRKuznetsov}
A.~Kuznetsov, \emph{Hochschild homology and semiorthogonal decompositions},
\href{https://arxiv.org/abs/0904.4330}{arXiv:0904.4330}.
射影核を使うadmissible部分圏のHochschild理論。

\bibitem{GRStacks}
The Stacks Project Authors, \emph{The Stacks Project},
\href{https://stacks.math.columbia.edu/tag/0B91}{Tag 0B91}, Lemma~36.30.4:
flat proper有限表示射のperfect直接像と任意の導来底変更。

\bibitem{GRBFK}
M.~Ballard, D.~Favero and L.~Katzarkov,
\emph{A category of kernels for equivariant factorizations and its implications for Hodge theory},
Publ. Math. Inst. Hautes \`Etudes Sci. \textbf{120} (2014), 1--111.
\href{https://doi.org/10.1007/s10240-013-0059-9}{doi:10.1007/s10240-013-0059-9}, Theorem~5.39.

\bibitem{GRLZ}
X.~Lin and S.~Zhang, \emph{Serre algebra, matrix factorization and categorical Torelli theorem for hypersurfaces},
\href{https://arxiv.org/abs/2310.09927v2}{arXiv:2310.09927v2} (2024),
Theorem~3.5, Remark~3.6 and \S5. 本節で引用するのは体上の計算。

\bibitem{GRLZ26}
X.~Lin and S.~Zhang, \emph{Categorical Torelli and Fourier--Mukai partners of hypersurfaces},
\href{https://arxiv.org/abs/2611.05915v1}{arXiv:2611.05915v1} (2026-11-04), \S2.1.
相対projected kernelの背景として参照し、本研究の非可換generic同値を与えるとはしない。

\bibitem{GRToenFinite}
B.~To\"en, \emph{Finitude homotopique des dg-alg\`ebres propres et lisses},
\href{https://doi.org/10.1112/plms/pdn027}{doi:10.1112/plms/pdn027}.
有限$A_n$決定性はsmooth proper性を仮定する。今回この定理を非可換族へ適用してはいない。

\bibitem{RSRBS} Jesse Burke and Greg Stevenson,
\emph{The derived category of a graded Gorenstein ring},
arXiv:1507.00830v1 (2015).
\url{https://arxiv.org/abs/1507.00830}.
参照：Definition 5.1, Proposition 4.18, Lemmas 6.1--6.2, Theorem 6.4.
\bibitem{RSRToen} Bertrand Toën,
\emph{Anneaux de définition des dg-algèbres propres et lisses},
arXiv:math/0611546v3 (2010).
\url{https://arxiv.org/abs/math/0611546}.
参照：Théorème 1, F4, F6.
\bibitem{RSRTV} Bertrand Toën and Michel Vaquié,
\emph{Moduli of objects in dg-categories},
arXiv:math/0503269v5 (2007).
\url{https://arxiv.org/abs/math/0503269}.
参照：Lemmas 2.8, 2.10, Corollary 2.13, Theorem 3.6, Corollaries 3.15, 3.17, 3.19.
\bibitem{RSROrlov} Dmitri Orlov,
\emph{Smooth and proper noncommutative schemes and gluing of DG categories},
arXiv:1402.7364v5 (2015).
\url{https://arxiv.org/abs/1402.7364}.
参照：Proposition 3.8, Theorem 3.25（体上）.
\bibitem{RSRRepo} nc-cubic project,
\emph{一般次元の局所代数的模型、非楕円型量子化、およびrotation付き可換復元},
commit \texttt{4f2eb328a2249a173f0c658cdb744277f041ff1d}, 2026年11月10日JST.
\url{https://github.com/uedakazushi/nc-cubic}.
一次原稿：\texttt{papers/sections/local\_projective\_model.tex},
\texttt{nonelliptic\_quantization.tex}, \texttt{rotation\_models.tex}.
局所模型の入力、親原稿の形式的復元と、本稿の双加群による局所実現証明を区別する。



\bibitem{P2Kontsevich}
M.~Kontsevich,
Deformation quantization of Poisson manifolds,
Lett. Math. Phys. \textbf{66} (2003), 157--216.
\bibitem{P2LPT}
F.~Loray, J.~V.~Pereira and F.~Touzet,
Foliations with trivial canonical bundle on Fano 3-folds,
Math. Nachr. \textbf{286} (2013), 921--940.
\bibitem{P2Mayanskiy}
E.~Mayanskiy,
Poisson cohomology of two Fano threefolds,
\url{https://arxiv.org/abs/1303.5916}.
